\documentclass[12pt,a4paper]{article}

\usepackage[utf8]{inputenc}
\usepackage[T1]{fontenc}
\usepackage[spanish]{babel}
\usepackage{geometry}
\geometry{margin=2.2cm}
\usepackage{amsmath,amssymb}
\usepackage{tikz}
\usetikzlibrary{shapes.geometric, arrows.meta, positioning, fit, backgrounds, calc}
\usepackage{booktabs}
\usepackage{hyperref}
\usepackage{fancyhdr}
\usepackage{titlesec}
\usepackage{enumitem}
\usepackage{xcolor}
\usepackage{graphicx}
\usepackage{float}
\usepackage{listings}
\usepackage{xcolor}

\definecolor{codegreen}{rgb}{0,0.6,0}
\definecolor{codegray}{rgb}{0.5,0.5,0.5}
\definecolor{codepurple}{rgb}{0.58,0,0.82}
\definecolor{backcolour}{rgb}{0.95,0.95,0.92}
\lstdefinestyle{mystyle}{
    backgroundcolor=\color{backcolour},   
    commentstyle=\color{codegreen},
    keywordstyle=\color{blue}\bfseries,
    numberstyle=\tiny\color{codegray},
    stringstyle=\color{codepurple},
    basicstyle=\ttfamily\footnotesize,
    breakatwhitespace=false,         
    breaklines=true,                 
    captionpos=b,                    
    keepspaces=true,                 
    numbers=left,                    
    numbersep=5pt,                  
    showspaces=false,                
    showstringspaces=false,
    showtabs=false,                  
    tabsize=2
}
\lstset{style=mystyle}

\hypersetup{
    colorlinks=true,
    linkcolor=blue!70!black,
    urlcolor=blue!80!black
}

\pagestyle{fancy}
\fancyhf{}
\setlength{\headheight}{15pt}
\fancyhead[L]{Ingeniería de Software -- Simulador Multirotor}
\fancyhead[R]{\thepage}
\renewcommand{\headrulewidth}{0.4pt}

\titleformat{\section}{\Large\bfseries\color{blue!70!black}}{\thesection}{1em}{}
\titleformat{\subsection}{\large\bfseries\color{blue!50!black}}{\thesubsection}{1em}{}

\title{\textbf{Informe de Ingeniería de Software}\\[0.3cm]
\large Arquitectura, Implementación y Diseño del Simulador Multirotor}
\author{Análisis del repositorio \texttt{github.com/ChristianGonper/modelado-control-dron}}
\date{\today}

\begin{document}

\maketitle

\begin{abstract}
Este informe analiza el repositorio desde la perspectiva de la \textbf{Ingeniería de Software}. Se describe la arquitectura modular, los contratos de datos, la implementación de los modelos físicos, el sistema de control (clásico y neuronal), 
el pipeline de entrenamiento, la orquestación y las decisiones de diseño que permiten la extensibilidad y la reproducibilidad del simulador.
\end{abstract}

\newpage
\tableofcontents
\newpage

\section{Visión General de la Arquitectura}

El simulador sigue una arquitectura \textbf{en capas} con alta separación de responsabilidades:

\begin{center}
\begin{tikzpicture}[
    node distance=0.8cm,
    box/.style={draw, rounded corners, minimum width=8cm, minimum height=1cm, align=center, font=\small},
    layer1/.style={box, fill=blue!10},
    layer2/.style={box, fill=green!10},
    layer3/.style={box, fill=orange!10},
    layer4/.style={box, fill=purple!10},
]
\node[layer1] (app) {Application Layer -- \texttt{app.py}, \texttt{runner.py}, CLI};
\node[layer2, below=of app] (control) {Control Layer -- \texttt{control/} (PD, MLP, GRU, LSTM)};
\node[layer3, below=of control] (dynamics) {Dynamics \& Physics Layer -- \texttt{dynamics/}, \texttt{core/}};
\node[layer4, below=of dynamics] (data) {Data \& Dataset Layer -- \texttt{dataset/}, telemetry, scenarios};
\end{tikzpicture}
\end{center}

\textbf{Principios clave}:
\begin{itemize}
    \item \textbf{Contrato primero}: todo se comunica a través de dataclasses inmutables (\texttt{VehicleState}, \texttt{VehicleIntent}, \texttt{VehicleCommand}).
    \item \textbf{Intercambiabilidad}: cualquier controlador que implemente \texttt{ControllerContract} puede usarse sin cambiar el runner.
    \item \textbf{Reproducibilidad}: semillas, normalización train-only, y logging completo.
\end{itemize}

\section{Estructura de Paquetes y Módulos}

\begin{lstlisting}[language=Python, caption=Estructura principal del paquete]
src/simulador_multirotor/
|-- __init__.py
|-- app.py                  # CLI y punto de entrada
|-- runner.py               # Orquestador principal
|-- benchmark.py
|-- control/
|   |-- __init__.py
|   |-- contract.py         # ControllerContract (Protocol)
|   |-- cascade.py          # Control PD en cascada
|   |-- mlp.py              # MLP + entrenamiento
|   `-- recurrent.py        # GRU/LSTM + entrenamiento
|-- core/
|   |-- contracts.py        # VehicleState, RotorGeometry, VehicleIntent...
|   `-- attitude.py         # Cuaterniones y helpers matematicos
|-- dynamics/
|   |-- rigid_body.py       # 6DOF + actuadores + mixer
|   `-- aerodynamics.py     # Drag, viento OU
|-- dataset/                # Ventanas, splits, features
|-- telemetry/
|-- trajectories/
|-- scenarios/
|-- metrics/
`-- visualization/
\end{lstlisting}

\section{Contratos de Datos (core/contracts.py)}

El corazón del diseño son las \textbf{dataclasses inmutables} con validación estricta en \texttt{\_\_post\_init\_\_}:

\begin{itemize}
    \item \texttt{VehicleState}: estado físico completo (posición, cuaternión, velocidades).
    \item \texttt{RotorGeometry}: parámetros físicos de cada rotor.
    \item \texttt{VehicleIntent}: alto nivel (empuje colectivo + torque cuerpo).
    \item \texttt{VehicleCommand}: lo que se aplica a la planta.
    \item \texttt{TrajectoryReference}: referencia de trayectoria.
\end{itemize}

\textbf{Ventaja}: inmutabilidad + validación temprana evita errores en tiempo de ejecución y facilita el razonamiento.

\section{Implementación de la Física}

\subsection{Matemáticas de Actitud (\texttt{core/attitude.py})}

Se implementan todas las operaciones de cuaterniones de forma vectorizada y numéricamente estable:

\begin{itemize}
    \item \texttt{quaternion\_multiply}
    \item \texttt{quaternion\_derivative} (usado en RK4)
    \item \texttt{quaternion\_error\_vector}
    \item \texttt{quaternion\_from\_thrust\_direction\_and\_yaw}
\end{itemize}

\subsection{Dinámica Rígida y Actuadores (\texttt{dynamics/rigid\_body.py})}

\begin{itemize}
    \item \texttt{RotorActuatorState.advance()} --- modelo de primer orden con saturación.
    \item \texttt{RotorMixer} --- pseudo-inversa + saturación por rotor.
    \item Integración RK4 del estado 6DOF.
\end{itemize}

\textbf{Nota de diseño}: el mixer y los actuadores están desacoplados del integrador, permitiendo cambiar la tasa de control independientemente de la física.

\section{Sistema de Control}

\subsection{Interfaz Común}

\texttt{ControllerContract} es un \texttt{Protocol} (duck typing estructural). Esto permite:

\begin{itemize}
    \item PD clásico (\texttt{cascade.py})
    \item MLP, GRU, LSTM
    \item Controladores externos o adaptados
\end{itemize}

Todos devuelven exactamente el mismo \texttt{VehicleCommand}.

\subsection{Control PD en Cascada}

Implementación limpia de dos bucles con feedforward y saturación por eje. Las ganancias están parametrizadas y validadas.

\subsection{Controladores Neuronales}

\begin{itemize}
    \item \texttt{MLPTrainingConfig} + \texttt{MLPArchitectureConfig}
    \item \texttt{RecurrentTrainingConfig}
    \item Guardado de \texttt{FeatureNormalizationStats} junto con los pesos.
    \item Inferencia idéntica a entrenamiento gracias al checkpoint completo.
\end{itemize}

\section{Pipeline de Datos y Entrenamiento}

\begin{center}
\resizebox{\textwidth}{!}{
\begin{tikzpicture}[
    node distance=0.6cm,
    box/.style={draw, rounded corners, minimum width=2.8cm, minimum height=0.9cm, align=center, font=\footnotesize},
    arrow/.style={-{Stealth[length=2mm]}}
]
\node[box, fill=blue!15] (raw) {Raw Telemetry};
\node[box, fill=green!15, right=of raw] (episodes) {Episodes};
\node[box, fill=orange!15, right=of episodes] (windows) {Windows};
\node[box, fill=purple!15, right=of windows] (train) {Train / Val Split};
\node[box, fill=red!15, right=of train] (model) {Model Training};
\node[box, fill=teal!15, right=of model] (ckpt) {Checkpoint};

\draw[arrow] (raw) -- (episodes);
\draw[arrow] (episodes) -- (windows);
\draw[arrow] (windows) -- (train);
\draw[arrow] (train) -- (model);
\draw[arrow] (model) -- (ckpt);
\end{tikzpicture}
}
\end{center}

\textbf{Características clave}:
\begin{itemize}
    \item Ventanas deslizantes con \texttt{stride}.
    \item Dos modos de features: \texttt{raw\_observation} y \texttt{observation\_plus\_tracking\_errors}.
    \item Normalización \textbf{solo con estadísticas de train}.
    \item Reproducibilidad total con semillas.
\end{itemize}

\section{Orquestación y Runner}

El \texttt{SimulationRunner} es el componente central:

\begin{itemize}
    \item Separa claramente \texttt{physics\_dt}, \texttt{control\_dt} y \texttt{telemetry\_dt}.
    \item Mantiene el estado de actuadores entre pasos de control.
    \item Registra metadatos completos (comando aplicado vs. comando deseado).
    \item Soporta escenarios externos JSON.
\end{itemize}

Esto permite experimentos de robustez y benchmark sin modificar controladores.

\section{Telemetría, Métricas y Escenarios}

\begin{itemize}
    \item \texttt{telemetry/export.py}: exporta historia completa a JSON.
    \item \texttt{metrics/report.py}: calcula RMSE de posición, velocidad, empuje, etc.
    \item \texttt{scenarios/}: generadores de escenarios mínimos y de referencia.
    \item \texttt{visualization/}: renderizado 2D/3D estático desde telemetría.
\end{itemize}

\section{Decisiones de Diseño y Patrones Aplicados}

\begin{enumerate}
    \item \textbf{Dataclasses frozen + slots}: inmutabilidad + bajo consumo de memoria.
    \item \textbf{Protocol-based design}: duck typing para controladores (sin herencia rígida).
    \item \textbf{Separación de tasas}: física, control y telemetría independientes.
    \item \textbf{Checkpointing rico}: arquitectura + normalización + pesos.
    \item \textbf{Reproducibilidad por defecto}: semillas en todos los componentes.
    \item \textbf{Validación temprana}: \texttt{\_\_post\_init\_\_} en todos los contratos.
\end{enumerate}

\section{Diagrama de Clases Principal (Resumen)}

\begin{center}
\begin{tikzpicture}[
    node distance=0.5cm,
    class/.style={draw, rectangle, minimum width=3.2cm, minimum height=2.8cm, align=left, font=\scriptsize},
    arrow/.style={-{Stealth[length=2mm]}, thick}
]
% Core
\node[class, fill=blue!8] (state) at (0,0) {
    \textbf{VehicleState}\\
    position\_m\\
    orientation\_wxyz\\
    linear\_velocity\_m\_s\\
    angular\_velocity\_rad\_s
};

\node[class, fill=green!8] (intent) at (5,0) {
    \textbf{VehicleIntent}\\
    collective\_thrust\_newton\\
    body\_torque\_nm
};

\node[class, fill=orange!8] (command) at (10,0) {
    \textbf{VehicleCommand}\\
    intent: VehicleIntent
};

% Control
\node[class, fill=purple!8] (controller) at (5,-4) {
    \textbf{ControllerContract}\\
    +compute\_action()\\
    (Protocol)
};

\node[class, fill=red!8] (cascade) at (0,-4) {
    \textbf{CascadeController}\\
    PositionLoop\\
    AttitudeLoop
};

\node[class, fill=teal!8] (mlp) at (10,-4) {
    \textbf{MLPController}\\
    model: nn.Module\\
    norm\_stats
};

\draw[arrow] (controller) -- (cascade);
\draw[arrow] (controller) -- (mlp);
\draw[arrow] (state) -- (controller);
\draw[arrow] (intent) -- (command);
\end{tikzpicture}
\end{center}

\section{Conclusiones de Ingeniería de Software}

El repositorio demuestra una \textbf{excelente práctica de ingeniería de software} para un proyecto de investigación:

\begin{itemize}
    \item Alta cohesión y bajo acoplamiento gracias a los contratos inmutables.
    \item Extensibilidad real: añadir un nuevo controlador neuronal es cuestión de implementar \texttt{compute\_action}.
    \item Reproducibilidad y trazabilidad completas.
    \item Separación clara entre física y software que facilita tanto el mantenimiento como la experimentación.
\end{itemize}

Este diseño permite pasar de un controlador PD clásico a modelos neuronales (MLP/GRU/LSTM) sin tocar ni una línea del runner ni de la planta física.

\vspace{1cm}
\begin{center}
\textit{Informe generado a partir del análisis del repositorio}\\
\texttt{https://github.com/ChristianGonper/modelado-control-dron}\\
\textit{Fecha: \today}
\end{center}

\end{document}